% Generated draft: requires mathematical and source review.
\documentclass[11pt]{article}
\usepackage[margin=25mm]{geometry}
\usepackage{fontspec}
\setmainfont{Latin Modern Roman}
\usepackage{amsmath,amssymb,mathtools,unicode-math}
\setmathfont{Latin Modern Math}
\usepackage{enumitem,needspace,xurl,hyperref}
\setlength{\parindent}{0pt}
\setlength{\parskip}{5pt}
\setlength{\emergencystretch}{3em}
\newcommand{\N}{\mathbb{N}}
\newcommand{\Z}{\mathbb{Z}}
\newcommand{\Q}{\mathbb{Q}}
\newcommand{\R}{\mathbb{R}}
\newcommand{\C}{\mathbb{C}}
\newcommand{\F}{\mathbb{F}}
\newcommand{\A}{\mathbb{A}}
\newcommand{\PP}{\mathbb{P}}
\newcommand{\E}{\mathbb{E}}
\newcommand{\eps}{\varepsilon}
\newcommand{\dd}{\,\mathrm{d}}
\begin{document}
\section*{UnsolvedMath ID: 6800026}
\textbf{Source problem code:} catalogue entry
\subsection{Definitions and mathematical statement}
In the K\"ahler setting indicated by the title, let $\omega(t)$ be a normalized K\"ahler--Ricci flow. A common normalization has
\[
\frac{\partial\omega}{\partial t}=-\operatorname{Ric}(\omega)+\lambda\omega,
\]
with the constant and initial K\"ahler class fixed by the geometric setting. For example, on a Fano manifold with $[\omega(0)]=2\pi c_1(M)$, one uses $\lambda=1$. This is the context highlighted by the supplied summary, not an explicit hypothesis of the extracted question.

Does the flow converge, as $t\to\infty$, in Gromov--Hausdorff sense to a generalized K\"ahler--Einstein space? A metric convergence assertion seeks a compact metric space $(Z,d_Z)$ such that
\[
\lim_{t\to\infty}d_{GH}((M,d_{g(t)}),(Z,d_Z))=0.
\]
Here $d_{GH}$ is the infimum of Hausdorff distances between isometric copies in common metric spaces. The requested convergence concerns the whole flow, not just a convergent subsequence.

The initial manifold category, flow normalization, long-time existence hypotheses, and meaning of generalized K\"ahler--Einstein are absent. On a smooth regular locus, an ordinary K\"ahler--Einstein condition is $\operatorname{Ric}(\omega)=\lambda\omega$. A K\"ahler--Ricci soliton instead allows a Lie-derivative term, $\operatorname{Ric}(\omega)-\lambda\omega=\mathcal L_V\omega$ under one sign convention. The summary mentions solitons and possibly singular $\mathbb Q$-Fano limits; these are not automatically identical to a specified generalized Einstein target. A precise theorem must fix that distinction and the singular-space category.
\subsection{Short English statement}
Does normalized K\"ahler--Ricci flow have a single Gromov--Hausdorff limit with the intended generalized K\"ahler--Einstein structure? The hypotheses, normalization, and distinction between Einstein and soliton limits must be specified.
\subsection{Sources}
\begin{itemize}
\item Morgan and Pansu, \emph{A List of Open Problems in Differential Geometry} (2018), Question environment 26, ``Gromov--Hausdorff convergence of K\"ahler Ricci flow,'' proposed by Gang Tian, BICMR, Peking University, \url{https://www.imo.universite-paris-saclay.fr/~pansu/problems_MTDG.tex}; accessed 2026-07-29, source-TeX extraction.
\item The supplied summary mentions Tian's program, Hamilton--Tian--Zhang convergence work, Chen--Donaldson--Sun, and singular/soliton limits without complete publication identifiers or exact theorem hypotheses.
\item Local registry \texttt{data/unsolvedmath/problems.json}, permanent UnsolvedMath ID $6800026$, source code catalogue entry, same titled problem. Dataset attribution: Ulam AI and the UnsolvedMath Contributors.
\end{itemize}
\end{document}
